% Section 7: Conjectures and open problems.
\section{Conjectures and open problems}\label{sec:conj}

The theorems of Sections~\ref{sec:equilibrium}--\ref{sec:natural} concern the lattice, apart from the potential theory of
Section~\ref{sec:equilibrium}, which holds for general densities. This section states what they and the evidence of
Section~\ref{sec:examples} suggest for other node sets, and what is missing for proofs.

\subsection{The natural-index law}\label{sec:conj-law}
Call $\Lambda$ \emph{regular} if $\lambda_m^{1/p}$ has an asymptotic expansion $m+a-1+\sum_{i\ge1}\gamma_im^{-r_i}$ with
exponents $0<r_1<r_2<\dotsb\to\infty$, or if $\Lambda$ differs from such a set in finitely many elements. Then $Z_\Lambda$ is
meromorphic, $Z_\Lambda(0)$ and $Z_\Lambda(-j)$ are defined by continuation, and $w\mapsto Z_\Lambda(-w)-\zeta(-pw,a)$ has
poles only at points $q$ produced by the corrections $\gamma_im^{-r_i}$ (at $q=r_i/p$ to first order).

\begin{conjecture}[natural-index law]\label{conj:natural}
Let $\Lambda$ be regular and $N_\Lambda:=n+\frac\alpha2-Z_\Lambda(0)$. Then $n^pb_n/c=(n/N_\Lambda)^pF$, where $F$ has an
asymptotic expansion in powers of $1/N_\Lambda$ and of the spectral monomials below, with
\begin{enumerate}
\item[(a)] no term $N_\Lambda^{-1}$, and the term $f_2N_\Lambda^{-2}$, with $f_2$ as in~\eqref{eq:f2};
\item[(b)] for $j\ge1$ the term $pW(j)Z_\Lambda(-j)N_\Lambda^{-1-jp}$;
\item[(c)] for every pole $q\ne-1/p$ of $Z_\Lambda(-w)-\zeta(-pw)$, with residue $R_q$, the term
$pR_qW(q)\pi\cot(\pi q)N_\Lambda^{-1-pq}$, with the value $-(4c)^q\Lt(q)$ at $q\in-\frac12-\N_0$.
\end{enumerate}
Coefficients of coincident exponents add. The statement concerns the leading power of each datum in~(b) and~(c); it
says nothing about the terms of order $N_\Lambda^{-2-pq}$ (Remark~\ref{rem:sublead}).
\end{conjecture}

For the lattice, (a) and (b) are Theorem~\ref{thm:N6} (with $Z_\Lambda(0)=-\frac12$, $W(j)p\zeta(-jp)=\varepsilon_j$), and (c)
is empty. For Bessel zeros the whole expansion is Proposition~\ref{prop:bessel}. For Hurwitz lattices, synthetic
corrections, smoothed drums and the oscillators $q^N$ the evidence is in Section~\ref{sec:examples}. In the form~(b)--(c) the
two kinds of spectral data are the residues of one integrand,
\[
  \frac pn\cdot\frac1{2\pi i}\int_{\uparrow}-W(w)\,\pi\cot(\pi w)\,Z_\Lambda(-w)\,n^{-wp}\diff w ,
\]
at the integers (Theorem~\ref{thm:K}) and at the poles of $Z_\Lambda$ (Proposition~\ref{prop:M12}). The Weyl pole of
$Z_\Lambda$ is excluded: it is the leading term $n$ of the soft-edge condition, not a correction. A moved Weyl pole, as for
$\lambda_k=k^p+\kappa$, is included (\S\ref{sec:num-pole}). If a pole meets an integer, the two residues diverge separately
and their sum produces a term $n^{-1-jp}\log n$; one weak test at $p=m=\frac12$ is consistent with the predicted coefficient.

Two consequences make the law a spectral dictionary. Through $Z_\Lambda(0)$, the $1/n$ coefficient $d_1=-p\alpha/2+pZ_\Lambda(0)$
reads off the index offset of the spectrum, such as the constant $a$ of a Hurwitz lattice (\S\ref{sec:hurwitz}) or the
corner and curvature constant of a drum (\S\ref{sec:drums}). Through $Z_\Lambda(-j)$ and the poles, it reads off the
coefficients of the heat trace of the operator whose spectrum is $\Lambda$, weighted by $W(q)\pi\cot(\pi q)$.

\begin{remark}[subleading pole terms]\label{rem:sublead}
Each pole also contributes at order $n^{-2-pq}$, one order below its leading term, and this contribution is not removed
by $\cot(\pi q)=0$. For Bessel zeros at $p=2$ the McMahon correction has $q=\frac12$; its leading term vanishes, and its
subleading term at $n^{-3}$ cancels the zeta term $2pcZ_\Lambda(-1)$ exactly, which is why Proposition~\ref{prop:bessel} is a
pure index shift; a pure Hurwitz lattice with the same offset does carry the zeta term at $n^{-3}$. For the synthetic
family the subleading term is~\eqref{eq:S4} when $q_1\le1$:
\end{remark}

\begin{conjecture}\label{conj:S4}
For $\lambda_k=(k+\gamma k^{-m})^p$, $m>0$, $p\ne1$, $q_1=(m+1)/p\le1$, the part linear in $\gamma$ of the $n^{-2-m}$ coefficient
of $n^pb_n/c$ is $\gamma S$ with $S$ as in~\eqref{eq:S4}.
\end{conjecture}

Both terms of~\eqref{eq:S4} are of the type that Conjecture~\ref{conj:natural} does not describe: a change of node density
on the band, and an index shift of a first-order term. In Theorem~\ref{thm:N6} the same shift is visible as the factor
$(n/N)^p$ in front of every layer. This suggests that every first-order term is to be evaluated at the natural index,
plus a density term.

\subsection{The natural form to all orders}\label{sec:conj-even}
\begin{conjecture}\label{conj:S1Nprime}
For the lattice and $p\ne1$, $F:=(N/n)^pn^pb_n/c$ depends on $n$ only through $N^{-2}$ and the monomials
$X_j:=\varepsilon_jN^{-1-jp}$: $F=\sum_\beta c_\beta(N^{-2})\prod_jX_j^{\beta_j}$, with each $c_\beta$ a power series.
\end{conjecture}

Theorem~\ref{thm:N6} proves the instances $\beta=0$ to $N^{-6}$, $\beta=e_j$ to relative order $N^{-3}$ for every $j$ (and
$N^{-4}$ for $j=1$), and $\beta=2e_1$ to relative order $N^{-3}$, with $c_{2e_1}(0)=(3p+1)/(2p)$. The numbers
$(3p+1)/(2p)=1+(p+1)/(2p)$, like $\varepsilon_j$, do not depend on $\alpha$. A closed form for $c_0$ is known only at
$(p,\alpha)=(2,-\frac12)$, where it is $(1-x/4)^{-1/2}$ in $x=N^{-2}$; the hard-edge denominators $\prod_k((2k+1)p+2)$ of
Remark~\ref{rem:N6}(b) suggest a hypergeometric structure in $a'=\frac1p+\frac12$.

\subsection{Reflection and shift}\label{sec:conj-reflect}
The exponent $\alpha$ enters the jump data of Lemma~\ref{lem:N6a} in two ways: as the Bessel order $\nu$ in the coefficients
$a_k(\nu),b_k(\nu)$, which are polynomials in $\nu^2$, and as the Szeg\H{o} exponent $\beta$ in
$\rho^2=\phi^{2(\beta+1)}$ at the hard edge and in the factor $\phi^{-2(\beta+1)}$ at the soft edge, where
$\phi(z)=(z^{1/2}+(z-4c)^{1/2})/(2c^{1/2})$ is the exterior conformal map of $[0,4c]$. For the lattice $\nu=\beta=\alpha$.
Keeping $\nu$ and $\beta$ independent defines formal problems whose residue recursion (Section~\ref{sec:natural}) gives formal
series $b(n;\beta,\nu)$ in $1/n$.

\begin{theorem}[reflection]\label{thm:reflect}
As formal series, $b(n;\beta,\nu)=b(-n;-\beta-2,\nu)$; that is, $d_k(-\beta-2,\nu)=(-1)^kd_k(\beta,\nu)$ for all $k$.
\end{theorem}

\begin{proof}
With $T=\frac1{\sqrt2}\left(\begin{smallmatrix}1&1\\ i&-i\end{smallmatrix}\right)$ we have $\sigma_2T=T\sigma_3$, so
$P_0^\infty=T\mathfrak a_\infty^{\sigma_3}T^{-1}$ and $P_1=-c\sigma_2$ commute with $\sigma_2$. The core matrices $X_k$ and
$\sigma_3\Xi_k\sigma_3$ lie in the span of $I,\sigma_2$ for even $k$ and of $\sigma_3,i\sigma_1$ for odd $k$, so conjugation by
$\sigma_2$ multiplies them by $(-1)^k$, while $\sigma_2d^{\sigma_3}\sigma_2=d^{-\sigma_3}$ for scalars $d$. Hence
$\sigma_2\Delta_k(\beta)n^{-k}\sigma_2=\Delta_k(-\beta-2)(-n)^{-k}$, the recursion is covariant under conjugation by a constant
matrix, and $(Q_1)_{12}(Q_1)_{21}$ is invariant under conjugation by $\sigma_2$.
\end{proof}

\begin{lemma}[self-similar Buyarov--Rakhmanov identity]\label{lem:BR}
Let $\Phi$ denote the phase of the lattice equilibrium in the scaled frame, with hard-edge branch $2f^{1/2}$ and soft-edge
branch $\tilde\varphi$. Then $(1+pz\partial_z)\Phi=-2\log\phi$ up to the additive constants of the branches; in particular
$(1+pz\partial_z)(2f^{1/2})=2\operatorname{arsinh}(v/2)$ at $z=-v^2$ ($c=1$).
\end{lemma}

\begin{proof}
At the hard edge $2f^{1/2}=\pi\sqrt{-z}\,\mathfrak a(z)$, and the series of $\mathfrak a$ has the factors $1/(a'+k)$ with
$a'=\frac1p+\frac12$. Since $(1+pz\partial_z)[\sqrt{-z}\,z^k]=p(a'+k)\sqrt{-z}\,z^k$, these factors cancel term by term and
leave the arcsine series. The soft-edge branch is the analytic continuation of the same phase.
\end{proof}

Under the self-similar scaling of the lattice problem, $1+pz\partial_z$ is the derivative in $n$ of $n\Phi(xn^p)$; so
Lemma~\ref{lem:BR} says that the $n$-flow of the phase is generated by $\log\phi$, which also generates the $\beta$-flow. Exact
computer algebra (to order $n^{-4}$ symbolically in $(p,\nu,\beta)$) and $40$-digit computation (to order $n^{-8}$, and for
the zeta layer) support the following.

\begin{conjecture}[shift]\label{conj:shift}
$b(n;\beta,\nu)\,n^p/c=(n/N)^pF_p(N;\nu^2)$ with $N=n+\frac{\beta+1}2$ and $F_p$ independent of $\beta$; the same holds with
the band-factor parameters included.
\end{conjecture}

Theorem~\ref{thm:reflect} and Conjecture~\ref{conj:shift} together give, to all orders, the evenness of the integer layer in
$N$ (apply the reflection at $N'=n-\frac{\beta+1}2$) and in $\alpha$ (through $\nu^2$), and they identify the lattice's natural
index $n+\frac\alpha2+\frac12$ as the Szeg\H{o} shift. For varying weights $e^{-nV}$, evenness after a half-index shift is
proved through the string equations \cite{BleherIts}; an analogue of the string equation for the present formal problems
would be a natural route to Conjecture~\ref{conj:shift}.

\begin{problem}\label{prob:kernel2}
Find a second-order edge kernel: a function $g(w)$, analytic in $w$, with $g(j)=g_j$ for integers $j\ge1$, such that the
$N^{-2}$ dressing of a pole term at $q$ is $g(q)$. By Remark~\ref{rem:N6}(d) it should combine a polynomial in
$\theta=u\frac{\diff}{\diff u}$ acting on $(1-u)^{-1/2}$ with a multiple of $(1-u)^{1/2}$. The data of
Table~\ref{tab:subpole} and of drums are test cases.
\end{problem}

\subsection{What a proof needs}\label{sec:conj-proof}
The derivation of Conjecture~\ref{conj:natural}(c) from the analysis of \cite{PaperII} needs that analysis for a node set
other than the lattice, with three ingredients: the $g$-function of the constrained equilibrium of the $n$-dependent
continuum node density (Theorem~\ref{thm:Q}); the band factor given by the Cauchy integral of the jump of
$\log(1-e^{-2\pi it_\kappa})$ for the index function $t_\kappa(z)=\kappa(n^p/z)$ (Proposition~\ref{prop:pert}); and parametrix
corrections whose first-order dependence on the perturbation is $O(n^{-2-m})$. Under these, the first-order change of
$n^pb_n/c$ at order $n^{-1-m}$ is the edge shift of~\eqref{eq:Q} for the perturbed density, and
Proposition~\ref{prop:M12} gives its value. For Hurwitz lattices the band factor is unchanged in form (its moments are
Hurwitz zeta values), and what is missing is the index normalisation of the outer function, where the shift
$-p\delta$ of $d_1$ must arise.

\begin{remark}[a proved instance of the per-node law]\label{rem:S2}
Conjecture~\ref{conj:natural}, read at leading order per node, says that one node with $u=4c\lambda n^{-p}<1$ outside the band
contributes $-p(1-u)^{-1/2}/n$ to $y_n$ when it is removed: this is the derivative of the soft-edge condition of
Remark~\ref{rem:Q}(c). For one atom $k_0^{-p}$ of $\sigma_{p,s}$ this is proved, uniformly for $u_n=4c(k_0/n)^p\le1-\eta$, with
error $O(n^{-1-\epsilon})$, $\epsilon=\min(2p,1)$ for $p<1$ and $\epsilon=1$ for $p\ge1$, from the rank-one update of the Hankel
matrices and a contraction argument for the minimal solution of the recurrence at the deleted node; the proof uses only the
first-order expansions of $a_n$ and $b_n$.
% TODO: this is Theorem S2' of research/G1/RH/spectral/SPEC.md Sec. 3.2 (reviewed), not in any paper; outside the
% sources requested for this paper. Either include it with proof or cite it once it is written up.
\end{remark}

\subsection{Further problems}\label{sec:conj-further}
(a) \emph{The case $p=1$.} Drums and the measures $\sum_kk^{-s}\delta_{1/k}$ have $p=1$, where the two cotangents of
Theorem~\ref{thm:K} merge and \cite{PaperII} uses a different hard-edge analysis; Section~\ref{sec:drums} indicates that the
dictionary persists. (b) \emph{Higher layers.} The layers $\mu_1\mu_2$, $\mu_2^2$, $\mu_1^3$ control the coefficients below
$n^{-2-3p}$ and are not computed. (c) \emph{Irregular spectra.} For spectra with displacements of the order of the spacing
there is no deterministic expansion; Section~\ref{sec:noise} suggests that the variance of the irregular part is
$n\int\kappa(\omega)S_h(\omega)\diff\omega/\pi$ with a universal kernel spectrum $\kappa$. A closed form for $\kappa$, presumably from
the Airy kernel of the soft edge, is open, as is a proof that rigidity at wavelengths of two to forty spacings sets the
noise.
